\documentclass[conference]{IEEEtran}
\IEEEoverridecommandlockouts
\usepackage{cite}
\usepackage{amsmath,amssymb,amsfonts}
\usepackage{algorithmic}
\usepackage{graphicx}
\usepackage{textcomp}
\usepackage{xcolor}
\usepackage{booktabs}
\usepackage{multirow}
\usepackage{array}
\usepackage{placeins}
\usepackage{epsfig}
\usepackage{tikz}
\usetikzlibrary{positioning,arrows.meta,calc}

\def\BibTeX{{\rm B\kern-.05em{\sc i\kern-.025em b}\kern-.08em
    T\kern-.1667em\lower.7ex\hbox{E}\kern-.125emX}}

\begin{document}

\title{Edge AI-Based FPGA Accelerator for Spectrum Classification and Anomaly Detection}

\author{Author 1, Author 2, and Author 3
  %  \IEEEauthorblockN{Yash~Agrawal, Srijan~Mukhopadhyay, and Sandeep Joshi}
  %  \IEEEauthorblockA{Department of Electrical and Electronics Engineering\\ Birla Institute of Technology and Science, Pilani, Rajasthan 333031, India}
  %  \{f20230306, f20230334, sandeep.joshi\}@pilani.bits-pilani.ac.in
  }
\maketitle


\begin{abstract}
Radio Frequency (RF) spectrum monitoring is critical for modern electromagnetic signal environments, where non-cooperative signals, including jamming pulses, spoofing transmissions, and wide-band interference, must be classified and countered close to the sensor. This work presents a fully synthesizable field-programmable gate array (FPGA) accelerator for RF signal classification and anomaly detection at the sensor edge, without cloud connectivity or external memory. The system implements an eight-layer one-dimensional Residual Convolutional Neural Network (ResNet-1D) with 29,252 parameters. Hardware deployment uses Batch Normalization (BatchNorm) folding and integer-only 8-bit (INT8) inference with per-channel weight scales, 32-bit biases, and fixed-point requantization between layers, including scale-aligned residual additions. On the 30,000-sample RadioML 2018.01A test set, the integer model reaches 83.18\% accuracy against 83.26\% in 32-bit floating point (FP32), a loss of 0.08 percentage points. The hand-written Verilog accelerator time-multiplexes one $W{\times}W$ processing-element (PE) array over all layers under a descriptor-driven sequencer; the $W{=}8$ configuration closes timing at 175.4~MHz out of context and at 142.86~MHz inside a complete Zynq-7000 system on the ZedBoard, using 3,117 look-up tables (LUTs), 104 DSP slices and 14 block RAMs. One inference takes 20,315 cycles (142.2~$\mu$s). Measured on the board, the accelerator reproduces the golden model bit-exactly on all test vectors and completes an inference, including input transfer over AXI, in 161.9~$\mu$s, 60.5$\times$ faster than the same integer model running on the on-chip ARM Cortex-A9 processor (9.79~ms).
\end{abstract}

\begin{IEEEkeywords}
Edge AI, FPGA accelerator, RF anomaly detection, ResNet-1D, integer quantization, spectrum monitoring.
\end{IEEEkeywords}

\section{Introduction}

Modern electromagnetic spectrum monitoring systems must be designed to respond to the electromagnetic environment in near real-time. A jammer can suppress communications by barrage or spot-frequency jamming, a navigation spoofer can manipulate receiver state without the receiver's knowledge, and wideband interference from co-channel transmitters can degrade radar performance. The system must detect, classify, and respond to the anomalous waveform quickly and deterministically, close to the antenna.
The Field-Programmable Gate Arrays (FPGAs) present a compelling hardware platform. They synthesize completely pipelined custom datapaths for Convolutional
Neural Networks (CNNs) without any of the instruction fetch, decode or cache-miss latency experienced by Central
Processing Units (CPUs) and Graphics Processing Units (GPUs). This allows FPGA CNN inference accelerators to exploit spatial parallelism across channels and filter taps, and to realize deterministic latency for spectrum-monitoring applications.
However, the application of deep neural networks in FPGA logic imposes stringent engineering constraints: floating-point operations are prohibitively expensive in terms of look-up tables (LUTs) and digital signal processing (DSP) slices, per-channel running statistics in batch normalization (BN) operations require additional hardware, and vanishing gradients limit the training of deep plain CNNs.

In recent years, an end-to-end learning framework for automatic modulation classification based on raw in-phase (I) and quadrature (Q) samples has become the mainstream of automatic modulation classification (AMC) research \cite{b1, b2}.
The RadioML 2018.01A dataset \cite{b3} subsequently provided 24-class I/Q data at 26 Signal-to-Noise Ratio (SNR) levels and became the community standard benchmark. The authors in \cite{b4} demonstrated the efficiency of an attention-enhanced, transformer-based deep learning architecture. The authors in \cite{b5} have developed a systematic description of robust deep learning for physical layer communications.
Recent literature on FPGA-based AI accelerators \cite{b6} categorizes optimizations at the computational, memory, and system-level, emphasizing hand-designed RTL and extreme quantization as key for performance in resource-limited systems. Also, the authors in \cite{b7} proved that the models built using the attention mechanism obtain much better performance on wireless modulation classification.
Existing FPGA deep learning RF classifiers, such as the DeepRadio platform \cite{b8}, classify signals received through the RF front end to different modulation types in real time with low power. The RF System-on-Chip (RFSoC) work in \cite{b9} uses VGG10 and ResNet33 models on the RadioML 2018.01A dataset with FPGA deployment. Prior work in \cite{b10} demonstrated an FPGA-accelerated anomaly detector utilizing an unsupervised autoencoder for generic, out-of-distribution flagging; our work instead uses a supervised ResNet-1D that maps signals onto operational categories. In \cite{b11}, the authors achieved approximately 95\% classification accuracy using a ResNet with six residual stacks on the full 24-class RadioML dataset. Tools like hls4ml already implement optimizers that fuse batch normalizations following affine projection layers into a single node with equivalent fused weights \cite{b12}, and INT8 inference of ResNets on Xilinx devices is established practice. Recently, \cite{b13} surveyed deep transfer learning approaches for AMC, highlighting emerging architectures that improve generalization across diverse channel conditions. The authors in \cite{b14} demonstrated an FPGA-based real-time signal measurement system using deep learning accelerators.

These limitations motivate an end-to-end, hardware-accelerated pipeline in which every arithmetic step executed on the FPGA is specified, verified and measured. The major contributions of this work are summarized as follows.
\begin{itemize}
    \item A four-class operational signal grouping (Normal, Jammer, Spoofer, Interference) of the 24 RadioML 2018.01A modulations, used as a candidate task for demonstrating edge-deployed signal categorization.
    \item An integer-only INT8 inference scheme for the eight-layer ResNet-1D: per-channel symmetric weight quantization, 32-bit biases, fixed-point requantization between layers, and scale-aligned residual additions. Together with a bit-exact Python golden model, it keeps 83.18\% test accuracy (0.08 percentage points below FP32) while storing all parameters in under 30~KB of on-chip memory.
    \item A parameterizable, hand-written Verilog accelerator with no high-level synthesis (HLS): one $W{\times}W$ INT8 PE array, a descriptor-driven layer sequencer, three on-chip feature buffers and a six-stage requantization pipeline. A design-space exploration over $W\in\{4,8,16\}$ reports cycle counts and timing-closed frequencies; the $W{=}8$ design closes timing at 175.4~MHz out of context and at 142.86~MHz in a complete Zynq-7000 system, using 5.9\% of the LUTs of the xc7z020.
    \item On-board validation on the ZedBoard: the accelerator reproduces the golden model's logits bit-exactly on all test vectors and is 60.5$\times$ faster than the same integer model on the on-chip ARM Cortex-A9 processor.
\end{itemize}

\section{Proposed Architecture}
The RadioML 2018.01A dataset \cite{b3}, provides 2,555,904 labeled I/Q signal records across 24 modulation types at 26 SNR levels from $-20$~dB to $+30$~dB in 2~dB steps. Each record is a floating-point array of shape $(1024, 2)$ representing 1,024 complex-valued time samples with I/Q baseband components. The dataset models realistic channel impairments including additive white Gaussian noise (AWGN), multipath Rayleigh fading, carrier frequency offset, timing recovery errors, and hardware non-linearities such as IQ imbalance and phase noise.
Standard modulation recognition assigns one label per modulation type, yielding a 24-class problem with fine-grained distinctions, e.g., between 16 Quadrature Amplitude Modulation (QAM) and 32 QAM. For spectrum monitoring, the system often needs to know only whether a signal is benign or which operational category it belongs to. This work therefore groups the 24 modulations into four operational categories (Normal, Jammer, Spoofer, Interference) according to spectral characteristics and bandwidth occupancy. The grouping is a candidate mapping used to demonstrate edge-deployed signal categorization; validating it against operational electronic-warfare signal libraries is left to future work.
Table~\ref{tab:threat_mapping} lists the mapping.
\begin{table}[t]
\renewcommand{\arraystretch}{1.2}
\caption{Mapping of the 24 RadioML 2018.01A Modulations to the Four Categories}
\begin{center}
\footnotesize
\setlength{\tabcolsep}{3pt}
\begin{tabular}{c l p{5.3cm}}
\toprule
\textbf{ID} & \textbf{Category} & \textbf{RadioML modulations} \\
\midrule
0 & Normal & FM, GMSK, AM-DSB-WC, AM-DSB-SC, AM-SSB-WC, AM-SSB-SC \\
1 & Jammer & OOK, 64QAM, 128QAM, 256QAM \\
2 & Spoofer & BPSK, QPSK, 8PSK, 16PSK, 32PSK, OQPSK \\
3 & Interference & 16APSK, 32APSK, 64APSK, 128APSK, 16QAM, 32QAM, 4ASK, 8ASK \\
\bottomrule
\end{tabular}
\label{tab:threat_mapping}
\end{center}
\end{table}
For training and evaluation, records with SNR $\geq 0$~dB are kept, the first 128 time samples of each record are used and normalized by their maximum absolute value, each category is balanced to 50,000 samples, and the set is split 70/15/15 (stratified) into training, validation and a 30,000-sample test set.
% TODO(authors): Option A — insert the RadioML-based jamming / spoofing / interference dataset
% (7 classes, JSR sweep, energy-detector and expert-feature baselines) once the corrected Colab
% run (dataset v2) has completed. The accelerator supports it unchanged (7-class model already
% verified bit-exact in RTL).

A plain stacked 1D CNN applies convolutional layers sequentially with no bypass paths. During backpropagation, the gradient of the loss with respect to early-layer parameters is computed as a product of Jacobian matrices across every intermediate layer.
ResNet-1D maps a two-channel 128-point I/Q input to four class logits through eight computational stages. The architecture employs progressively increasing channel widths ($2 \rightarrow 16 \rightarrow 32 \rightarrow 64$) combined with stride-2 downsampling at two points to compress the temporal dimension from 128 to 32 timesteps before global average pooling. Table~\ref{tab:layerwise} provides a complete layer-wise breakdown.
%%%%%%%%%%%%%%%%%%%%%%%%%%%%%%%
\begin{table}[t]
\renewcommand{\arraystretch}{1.3}
\caption{ResNet-1D Layer-wise Architecture}
\begin{center}
\scriptsize
\setlength{\tabcolsep}{2.5pt}
\begin{tabular}{c l c c c r}
\toprule
\textbf{Layer index} & \textbf{Type} & \textbf{In (C$\times$L)} & \textbf{Out (C$\times$L)} & \textbf{k, s} & \textbf{Parameters} \\
\midrule
1 & Convolution+BN+ReLU & $2\!\times\!128$ & $16\!\times\!128$ & 7, 1 & 240 \\
2 & Convolution+BN+ReLU & $16\!\times\!128$ & $32\!\times\!64$ & 3, 2 & 1,568 \\
3--4 & ResBlk 1 (id) & $32\!\times\!64$ & $32\!\times\!64$ & 3, 1 & 9,280 \\
5--7 & ResBlk 2 (proj) & $32\!\times\!64$ & $64\!\times\!32$ & 3, 2 & 17,664 \\
8 & GAP & $64\!\times\!32$ & $64\!\times\!1$ & --- & 0 \\
9 & FC + Argmax & 64 & 4 & --- & 260 \\
\midrule
 & \multicolumn{4}{r}{\textbf{Total Trainable Parameters}} & \textbf{29,252} \\
\bottomrule
\end{tabular}
\label{tab:layerwise}
\end{center}
\footnotesize{ResBlk: Residual Block;
id: identity shortcut; proj: projection shortcut; GAP: Global Average Pooling;
FC: Fully Connected layer; ReLU: Rectified Linear Unit; C: channels; L: temporal length; s: stride; k: kernel size.}
\end{table}
%%%%%%%%%%%%%%%%%%%%%%%%%%%%%%%
The 1D convolution output for output channel index $c_o$ at time position $t$ is expressed as
\begin{equation}
y[c_o, t] = \sum_{c_i=0}^{C_{in}-1} \sum_{k=0}^{K-1} W[c_o, c_i, k] \, x[c_i, t \cdot s + k - p] + b[c_o],
\label{eq:conv1d}
\end{equation}
where $c_i$ is the input channel index, $K$ is the kernel size, $s$ the stride, $p$ the zero padding, $x$ is the input tensor, $W$ is the weight tensor, $[ \cdot, \cdot,\cdot]$ denotes tensor indexing, $C_{in}$ is the number of input channels, and $b[c_o]$ is the bias. Batch Normalization normalizes each channel, and the normalized output for each channel is expressed as
\begin{equation}
\hat{x}[c] = \gamma[c] \frac{x[c] - \mu[c]}{\sigma[c]} + \beta[c],
\label{eq:bn}
\end{equation}
where $x[c]$ is the input for channel $c$, $\mu[c]$ and $\sigma[c]$ are the channel's running mean and standard deviation, and $\gamma[c]$ and $\beta[c]$ are the learned scale and shift parameters, respectively. ReLU activation applies element-wise, $\text{ReLU}(x) = \max(0, x)$. The identity skip connection computes $y = \text{ReLU}(F(x) + x)$ when channel dimensions match; when they do not, a $1\times1$ projection convolution adjusts the shortcut as $y = \text{ReLU}(F(x) + W_p\ast x)$, where $F(x)$ is the residual mapping and $W_p$ is the $1\times1$ projection weight matrix. Global Average Pooling reduces each channel to a scalar, $ z[c] = \frac{1}{T} \sum_{t=0}^{T-1} x[c, t]$,
where $T$ is the temporal sequence length. Since $T = 32$, the hardware replaces the division by a rounding arithmetic right shift by 5 positions, $z[c] = (\sum_t x[c,t] + 16) \gg 5$, eliminating all division circuitry from the pooling datapath. The argmax classification selects $\hat{c} = \arg\max_{j} l[j]$ over the class logits $l[j]$, and an anomaly flag is asserted whenever $\hat{c} \neq 0$ (class 0 = Normal).

\section{INT8 Quantization and FPGA Implementation}

\subsection{Integer-Only Quantization}
Prior to quantization, every BatchNorm layer is algebraically absorbed into its preceding convolution. The folded weight is computed per output channel $c_o$ as
\begin{equation}
\tilde{W}[c_o, c_i, k] = \frac{\gamma[c_o]}{\sigma[c_o]} W[c_o, c_i, k],
\label{eq:bnfold_w}
\end{equation}
and the folded bias absorbs both the original convolution bias and the complete BN affine transformation,
\begin{equation}
\tilde{b}[c_o] = \frac{\gamma[c_o] \cdot (b[c_o] - \mu[c_o])}{\sigma[c_o]} + \beta[c_o].
\label{eq:bnfold_b}
\end{equation}
Applying this across all seven Conv--BN pairs (including the projection shortcut) removes every BN operation from the hardware.

This work employs symmetric post-training quantization (PTQ) \cite{b18} with integer-only arithmetic. Symmetric quantization (zero-point $z = 0$) keeps every MAC a single signed multiply without offset terms. Weights are quantized per output channel: the scale $s_W[c_o] = \max(|\tilde{W}[c_o,:,:]|)/127$ maps the largest weight of each channel to 127, and
\begin{equation}
W_q = \text{clamp}\left(\text{round}\left(\tilde{W}/s_W\right), -127, 127\right).
\label{eq:quant}
\end{equation}
Activations use one symmetric INT8 scale $s_x$ per tensor. The input scale is $1/127$ (I/Q samples normalized to $[-1, 1]$); the other activation scales are calibrated on 4,096 training frames at the 99.99th percentile of the absolute activation value. The bias is stored as a 32-bit integer at the accumulator scale, $b_q = \text{round}(\tilde{b}/(s_W s_x))$.

Each MAC multiplies a signed INT8 weight by a signed INT8 activation and accumulates into a signed 24-bit register. The worst-case layer has $C_{in}K = 64 \times 3 = 192$ terms, giving at most $192 \times 127 \times 128 = 3{,}121{,}152 < 2^{23}$, so the 24-bit accumulator cannot overflow. After the dot product, the result is rescaled to the next layer's INT8 grid by a fixed-point multiplier $M[c_o] = \text{round}(2^{SH} s_W[c_o] s_x / s_y)$ with $M < 2^{15}$ and a per-layer shift $SH$:
\begin{equation}
\begin{split}
v &= (acc + b_q)\,M + r_q\,M_s 2^{SK} + 2^{SH-1},\\
y_q &= \text{sat}_8\!\left(\text{ReLU}\!\left(v \gg SH\right)\right),
\end{split}
\label{eq:requant}
\end{equation}
where $r_q$ is the INT8 residual (shortcut) input, present only in the last convolution of each residual block, and $M_s 2^{SK} = \text{round}(2^{SH} s_r/s_y)$ aligns the shortcut's scale $s_r$ with the output scale $s_y$. This scale alignment is required for a correct residual addition, because the two operands of the addition carry different scales. The fully connected layer outputs raw 32-bit logits ($acc + b_q$) to the argmax. A Python reference implementation of exactly this arithmetic serves as the golden model for all hardware verification.

\subsection{Accelerator Architecture}
%%%%%%%%%%%%%%%%%%%%%%%%%%%%%%%%
\begin{figure}[t]
\centering
 \includegraphics[width=3.4in, height=1.8in]{impl.png}
\caption{Physical deployment on the ZedBoard and post-implementation floorplan of the complete Zynq-7000 system (PS7, AXI interconnect and the RadarNet-1D accelerator) on the xc7z020.}
\label{fig:impl}
\end{figure}
% TODO(authors): replace impl.png with the floorplan of the new implementation
% (Vivado: Open Implemented Design -> Device view) — the current image shows the previous design.
%%%%%%%%%%%%%%%%%%%%%%%%%%%%%%%%
\begin{figure*}[t]
\centering
\begin{minipage}[t]{0.49\textwidth}
  \centering
  \includegraphics[width=\textwidth]{radarnet_hw_overview.pdf}
  \caption{Top-level datapath of the RadarNet-1D accelerator: descriptor-driven sequencer, three feature buffers, weight and parameter ROMs, $W{\times}W$ PE array, accumulators, requantization pipeline with residual input, GAP unit and argmax.}
  \label{fig:hw_overview}
\end{minipage}\hfill
\begin{minipage}[t]{0.49\textwidth}
  \centering
  \includegraphics[width=\textwidth]{radarnet_dsp_mac_engine.pdf}
  \caption{One of the $W$ identical output lanes: $W$ DSP48 multipliers, a registered adder tree, the INT24 accumulator and the six-stage requantization pipeline of \eqref{eq:requant}; dotted lines mark pipeline registers ($W{=}8$ shown).}
  \label{fig:dsp_mac}
\end{minipage}
\end{figure*}
%%%%%%%%%%%%%%%%%%%%%%%%%%

Instead of instantiating one engine per layer, the accelerator time-multiplexes a single compute engine over all layers (Fig.~\ref{fig:hw_overview}). A sequencer reads a per-layer descriptor (channels, length, kernel, stride, padding, ReLU and residual flags, source/destination/shortcut buffers, ROM bases, and the requantization constants) and steps through the loop nest
output-channel group $o_g$ $\rightarrow$ time step $t$ $\rightarrow$ input-channel group $i_g$ $\rightarrow$ kernel tap $k$.
Every clock cycle it issues one word of $W$ input channels at one time position, together with a row of $W{\times}W$ weights, to the PE array, so $W^2$ MACs are performed per cycle and $W$ output channels are produced in parallel. Zero padding is applied in the operand register. The layer sequence is L1, L2, RB1C1, RB1C2 (+identity shortcut), RB2P (1$\times$1 projection), RB2C1, RB2C2 (+projection shortcut), GAP and FC; the projection shortcut and the fully connected layer run on the same engine as ordinary layers.

Activations are held in three ping-pong feature buffers whose words contain $W$ channels at one time position; one read therefore feeds the whole PE array, and the residual shortcut is read from a second buffer in the same cycle in which the corresponding convolution result is ready. Each feature map is overwritten as soon as no later layer needs it, so three 2-KB buffers suffice for the whole network. The weights are stored in a $462 \times 512$-bit ROM (one $W{\times}W$ block per row for $W{=}8$) and the per-channel bias and multiplier in a parameter ROM; all 28,640 INT8 weights, 308 INT32 biases and 304 multipliers (30.5~KB in total) reside in on-chip block RAM.

Fig.~\ref{fig:dsp_mac} details one output lane. The $W$ products are formed in DSP48 slices and reduced by a fully registered binary adder tree of $\log_2 W$ levels; the INT24 accumulator restarts at the first tap of each output and passes its result to the requantization pipeline R0--R5, which adds the bias, applies the multiplier $M$ in a DSP48, adds the scaled shortcut (a second DSP48), the rounding constant and the arithmetic shift, and finally applies ReLU and saturation before writing $W$ channels back to the destination buffer. GAP is computed by a lane-parallel unit, and the argmax compares the 32-bit FC logits sequentially. The datapath is pipelined throughout: two-stage address generation, BRAM reads with output registers, registered layer descriptors and precomputed loop bounds. Successive output groups enter the pipeline without bubbles, so the cycle count stays within 3\% of the ideal MAC-issue count.

For system integration, the accelerator is wrapped in an AXI4-Lite slave with control, status, sample, class, cycle-counter and logit registers, and is connected to the Zynq processing system (PS7) through its general-purpose AXI port; the programmable-logic clock is FCLK0 = 142.86~MHz. The ARM processor writes the 128 I/Q samples, starts an inference and reads back the class, the logits and the cycle count.

The complete co-design flow spans: (1) Python preprocessing; (2) PyTorch training, BN folding and integer calibration, which emits the weight and parameter ROM images and the layer-descriptor table; (3) a bit-exact Python golden model of the integer arithmetic; (4) RTL simulation in which the accelerator's four INT32 logits are compared with the golden model for every test vector; (5) Vivado implementation of the PS7 system; and (6) a bare-metal Vitis application that runs the same integer model in C on the ARM Cortex-A9, runs the accelerator through AXI, checks both against the golden logits and times both with the ARM global timer.

\section{Results}
This section evaluates the system along two axes: classification performance of the ResNet-1D model under the four-class grouping, and the resource, timing, latency and power characteristics of the accelerator, both from Vivado implementation and measured on the ZedBoard.
The ResNet-1D model was trained using PyTorch with the Adam optimizer \cite{b19} at an initial learning rate of $1 \times 10^{-3}$ and weight decay $1 \times 10^{-4}$. A cosine annealing learning rate schedule reduced the rate to $\eta_{\min} = 1 \times 10^{-6}$ over a 60-epoch budget. Early stopping with patience 10 terminated training at epoch 45. Cross-entropy loss was used with a mini-batch size of 256.

Table~\ref{tab:perclass} reports per-class precision, recall, and F1-score on the held-out test set of 30,000 balanced samples, for the FP32 model and for the integer (hardware) model. Normal signals reach an F1 above 99\%, owing to the distinct narrowband spectral characteristics of continuous-wave modulations, and Spoofer reaches 95\%, benefiting from the coherent phase structure of PSK modulations. Jammer and Interference reach about 69\%: the higher-order QAM modulations are split between these two groups, and the network cannot reliably separate QAM orders from 128 samples, which bounds the overall accuracy of this grouping.
%%%%%%%%%%%%%%%%%%%%%%%%%%
\begin{table}[t]
\renewcommand{\arraystretch}{1.3}
\caption{Per-Class Performance on the Test Set (30,000 Samples)}
\begin{center}
\footnotesize
\setlength{\tabcolsep}{2.6pt}
\begin{tabular}{lccccc}
\toprule
 & \multicolumn{3}{c}{\textbf{FP32}} & \textbf{INT8 HW} & \\
\textbf{Class} & \textbf{Prec. (\%)} & \textbf{Rec. (\%)} & \textbf{F1 (\%)} & \textbf{F1 (\%)} & \textbf{Support} \\
\midrule
Normal       & 99.22 & 99.45 & 99.33 & 99.37 & 7,500 \\
Jammer       & 70.45 & 68.11 & 69.26 & 68.93 & 7,500 \\
Spoofer      & 93.56 & 96.87 & 95.19 & 95.23 & 7,500 \\
Interference & 68.93 & 68.61 & 68.77 & 68.77 & 7,500 \\
\midrule
\textbf{Overall} & \textbf{83.04} & \textbf{83.26} & \textbf{83.14} & \textbf{83.07} & \textbf{30,000} \\
\bottomrule
\end{tabular}
\label{tab:perclass}
\end{center}
\footnotesize{Overall: macro-averaged precision and F1; recall equals accuracy for the balanced test set. INT8 HW: integer model executed bit-exactly by the accelerator.}
\end{table}
%%%%%%%%%%%%%%%%%%%%%%%%%%

Table~\ref{tab:quant} tracks test accuracy at each stage of the quantization pipeline. BN folding via \eqref{eq:bnfold_w} and \eqref{eq:bnfold_b} is algebraically exact and leaves the accuracy unchanged. The integer-only model, with INT8 weights and activations and the requantization of \eqref{eq:requant}, loses 0.08 percentage points and agrees with the FP32 prediction on 98.24\% of the test samples. Because the accelerator reproduces this integer model bit-exactly, its accuracy is the hardware accuracy.
%%%%%%%%%%%%%%%%%%%%%%%%%%%%%
\begin{table}[t]
\renewcommand{\arraystretch}{1.3}
\caption{Test Accuracy at Each Stage of the Quantization Pipeline}
\begin{center}
\begin{tabular}{lcc}
\toprule
\textbf{Pipeline Stage} & \textbf{Accuracy (\%)} & \textbf{$\Delta$ (pp)} \\
\midrule
FP32 Baseline  & 83.26 & --- \\
After BN Folding (FP32) & 83.26 & 0.00 \\
Integer-only INT8 (hardware arithmetic) & 83.18 & $-$0.08 \\
\bottomrule
\end{tabular}
\label{tab:quant}
\end{center}
\end{table}
%%%%%%%%%%%%%%%%%%%%%%%%%%%%%%%%%%%

%%%%%%%%%%%%%%%%%%%%%%%%%%%%%%
\begin{table*}[t]
\renewcommand{\arraystretch}{1.3}
\caption{Comparison with FPGA-Based RF Classifiers}
\begin{center}
\begin{tabular}{lccccc}
\toprule
\textbf{Metric} & \textbf{Proposed} & \textbf{DeepRadio \cite{b8}} & \textbf{hls4ml \cite{b12}} & \textbf{FINN-based \cite{b20}} & \textbf{RFSoC-AMC \cite{b9}} \\
\midrule
\textbf{Model Architecture} & ResNet-1D (8-layer) & Custom CNN & Configurable & BNN & VGG10/ResNet33 \\
\textbf{Dataset} & RadioML 2018.01A & Proprietary RF & RadioML & RadioML & RadioML 2018.01A \\
\textbf{Task} & 4-class grouping & Per-modulation & Per-modulation & Per-modulation & Per-modulation \\
\textbf{Quantization} & INT8, integer-only & Fixed-point 16-bit & Mixed fixed-point & 1-bit (Binary) & INT8 (HLS) \\
\textbf{HLS Dependency} & None (custom RTL) & Partial & Full HLS & Full HLS & Full HLS \\
\textbf{Accuracy (\%)} & 83.18 & $\sim$80 & $\sim$85 & $\sim$80.1 & $\sim$90 \\
\textbf{Latency ($\mu$s)} & 142$^{a}$ / 162$^{b}$ & $\sim$5 & $\sim$1--10 & $\sim$283 & $\sim$10 \\
\textbf{Power (W)} & 0.29 (accelerator, est.) & $\sim$2--5 & $\sim$1--3 & $\sim$3.86 & $\sim$5--10 \\
\textbf{Energy/Inf. ($\mu$J)} & $\approx$41 (est.) & $\sim$10--25 & $\sim$1--30 & $\sim$1.21 & $\sim$50--100 \\
\textbf{LUT / DSP} & 3,117 / 104 (whole PL) & --- & --- & --- & --- \\
\textbf{Parameter Storage (KB)} & 30.5 & $>$100 & Variable & $\sim$187.5 & $>$500 \\
\textbf{Target FPGA} & Zynq-7000 & Zynq UltraScale+ & Various & Zynq-7000 & RFSoC ZU+ \\
\textbf{Parameters} & 29,252 & $>$100K & Variable & $>$1.5M & $>$1M \\
\textbf{On-board validation} & Bit-exact, measured & --- & --- & --- & --- \\
\bottomrule
\end{tabular}
\label{tab:comparison}
\end{center}
\footnotesize{BNN: Binary Neural Network. est.: Vivado vectorless estimate. $^{a}$compute, 142.86~MHz. $^{b}$measured on board, end to end including AXI input transfer.}
\end{table*}
%%%%%%%%%%%%%%%%%%%%%%%%%%%%%%
\subsection{Design-Space Exploration and Timing}
Table~\ref{tab:dse} compares three PE-array sizes, each implemented out of context on the xc7z020 and re-implemented at the frequency reported, with the clock modelled on a global buffer; a configuration is counted only if setup and hold slack of all internal register-to-register paths are non-negative. The cycle count is measured from start to result in RTL simulation, including the 128-cycle frame load. MAC utilization is the ratio of the 1,175,808 MACs of one inference to $W^2$ times the MAC-issue cycles. $W{=}8$ maps all 64 multipliers to DSP48 slices and gives the best balance; for $W{=}16$, the 256 multipliers exceed the 220 DSP slices of the device and are built from LUTs, which raises the LUT count roughly tenfold and lowers the clock.
%%%%%%%%%%%%%%%%%%%%%%%%%%%%%%%
\begin{table}[t]
\renewcommand{\arraystretch}{1.25}
\caption{Design-Space Exploration on xc7z020 (Out-of-Context)}
\begin{center}
\scriptsize
\setlength{\tabcolsep}{2.2pt}
\begin{tabular}{ccccccccc}
\toprule
$\boldsymbol{W}$ & \textbf{MAC util.} & \textbf{Cycles} & \textbf{$\boldsymbol{f_{\text{closed}}}$ (MHz)} & \textbf{Lat. ($\mu$s)} & \textbf{LUT} & \textbf{FF} & \textbf{BRAM} & \textbf{DSP} \\
\midrule
4  & 97.6\% & 76,123 & 169.5 & 449.1 & 1,551 & 1,197 & 11.5 & 28 \\
8  & 93.2\% & 20,315 & 175.4 & 115.8 & 2,671 & 1,948 & 14 & 104 \\
16 & 85.4\% & 5,855  & 138.9 & 42.2  & 24,792 & 14,523 & 41.5 & 32 \\
\bottomrule
\end{tabular}
\label{tab:dse}
\end{center}
\footnotesize{$W^2$ MACs per cycle. $f_{\text{closed}}$: frequency at which internal setup and hold slack are met. Lat.\ $=$ cycles$/f_{\text{closed}}$.}
\end{table}
%%%%%%%%%%%%%%%%%%%%%%%%%%%%%%%

The $W{=}8$ accelerator was integrated into the complete ZedBoard system (PS7, AXI interconnect and AXI4-Lite wrapper). At FCLK0 = 142.86~MHz, the routed system meets all timing constraints with a worst negative slack (WNS) of $+0.385$~ns, a worst hold slack of $+0.053$~ns and a worst pulse-width slack of $+2.520$~ns. The resource utilization of this system is listed in Table~\ref{tab:utilization} and its power estimate in Table~\ref{tab:power}. The same RTL also closes timing at 250~MHz with $W{=}16$ (all multipliers in DSP slices) inside a Zynq UltraScale+ ZU7EV system (WNS $+0.156$~ns), indicating that the architecture scales with the target device.

%%%%%%%%%%%%%%%%%%%%%%%%%%%%%%%
\begin{figure}[t]
\centering
\includegraphics[width=3.5in, height=1.7in]{behavorial_sim.png}
\caption{Behavioral simulation of the accelerator: frame load, layer sequence and result, with the class, the expected class of the golden model and the cycle count (20,315 cycles per inference).}
\label{fig:waveform}
\end{figure}
% TODO(authors): replace behavorial_sim.png with a waveform of the new design
% (Vivado xsim of tb_rn_top: result_valid, class_id, wv_expected, wv_hw_correct, cycle_count).
%%%%%%%%%%%%%%%%%%%%%%%%%%%%%%%

\begin{table}[t]
\renewcommand{\arraystretch}{1.3}
\caption{Resource Utilization of the ZedBoard System on xc7z020clg484-1 (Zynq-7000), Post-Placement}
\label{tab:utilization}
\centering
\begin{tabular}{lccc}
\toprule
\textbf{Resource} & \textbf{Used} & \textbf{Available} & \textbf{Utilization (\%)} \\
\midrule
LUT & 3,117 & 53,200 & 5.86 \\
FF & 2,448 & 106,400 & 2.30 \\
BRAM (36 Kb) & 14 & 140 & 10.00 \\
DSP48E1 & 104 & 220 & 47.27 \\
\bottomrule
\end{tabular}
\end{table}

\begin{table}[t]
\renewcommand{\arraystretch}{1.3}
\caption{On-Chip Power Estimate of the ZedBoard System at 142.86~MHz (Vivado, Vectorless)}
\label{tab:power}
\centering
\begin{tabular}{lc}
\toprule
\textbf{Component} & \textbf{Power (W)} \\
\midrule
\textbf{Total on-chip power} & \textbf{1.978} \\
\quad Processing system (PS7) & 1.530 \\
\quad Programmable logic, dynamic & 0.297 \\
\quad\quad of which RadarNet-1D accelerator & 0.290 \\
\quad\quad DSP / BRAM & 0.120 / 0.059 \\
\quad\quad Signals / logic / clocks & 0.062 / 0.034 / 0.022 \\
\quad Device static & 0.151 \\
\midrule
Junction temperature & 47.8$^{\circ}$C \\
\bottomrule
\end{tabular}
\end{table}

\subsection{Latency and On-Board Measurement}
Fig.~\ref{fig:waveform} shows the behavioral simulation. In RTL simulation, the four INT32 logits of the accelerator equal those of the golden model on all 50 test vectors, and one inference takes 20,315 cycles from start to result, including a 128-cycle frame load. At 142.86~MHz this corresponds to 142.2~$\mu$s.

The design was then run on the ZedBoard (Table~\ref{tab:board}). A bare-metal application executes the same integer model in C on the ARM Cortex-A9 (666.7~MHz, compiled with \texttt{-O3}) and on the accelerator, for 50 test vectors with 20 timed repetitions each, using the 333.3~MHz ARM global timer. Both produce logits identical to the golden model for all 50 vectors. On the board, the accelerator's cycle counter reads 23,013 cycles, because the 128 input samples arrive as individual AXI4-Lite writes from the processor rather than one per clock; including this transfer, one inference takes 161.9~$\mu$s, against 9.79~ms in software, a speed-up of 60.5$\times$. Based on the estimated accelerator power of 0.290~W, the accelerator consumes about 41~$\mu$J per inference (compute only); board-level power was not measured.

\begin{table}[t]
\renewcommand{\arraystretch}{1.3}
\caption{Measured on the ZedBoard: ARM Cortex-A9 Software vs.\ FPGA Accelerator (Same Integer Model, 50 Test Vectors)}
\label{tab:board}
\centering
\begin{tabular}{lcc}
\toprule
\textbf{Metric} & \textbf{ARM Cortex-A9} & \textbf{Accelerator} \\
\midrule
Logits bit-exact vs.\ golden model & 50 / 50 & 50 / 50 \\
Clock & 666.7~MHz & 142.86~MHz \\
Latency per inference & 9,787~$\mu$s & 161.9~$\mu$s$^{\dagger}$ \\
Cycle counter (start $\rightarrow$ result) & --- & 23,013 \\
Speed-up & 1$\times$ & 60.5$\times$ \\
\bottomrule
\end{tabular}

\vspace{2pt}
\footnotesize{$^{\dagger}$End to end, including the transfer of the 128 I/Q samples over AXI4-Lite; compute alone takes $\approx$20,190 cycles (141~$\mu$s).}
\end{table}

\subsection{Comparison with Prior FPGA RF Classifiers}
The proposed accelerator is contextualized relative to other FPGA-based RF classification accelerators in Table~\ref{tab:comparison}.
% TODO(authors): verify every value in the prior-work columns against the cited papers
% (same metric definitions, same dataset/class count) before submission; replace "~" estimates
% with the reported numbers or "n/r" (not reported).
The proposed accelerator targets a different operating point from throughput-oriented streaming designs: it runs a complete ResNet-1D with residual connections in 3,117 LUTs and 104 DSP slices of a low-cost Zynq-7000 device, keeps all 29,252 parameters on chip, and its integer arithmetic is verified bit-exactly from the Python model through RTL simulation to the physical board. Fully unrolled or streaming designs can reach lower latency at substantially higher resource cost; the $W$ parameter of the proposed architecture exposes this trade-off explicitly (Table~\ref{tab:dse}). Binarized approaches such as FINN reduce resources further but sacrifice representational capacity, while RFSoC-based deployments attain higher accuracy on the 24-class problem on substantially larger and more power-hungry devices.

\section{Conclusion}
This work has presented an end-to-end design of an RF signal classification and anomaly detection accelerator at the sensor edge. An eight-layer ResNet-1D with 29,252 parameters, trained on RadioML 2018.01A, reaches 83.26\% test accuracy in FP32 under a four-class operational grouping. BN folding and an integer-only INT8 scheme with per-channel weights, 32-bit biases, fixed-point requantization and scale-aligned residual additions reduce this by only 0.08 percentage points, to 83.18\%. A hand-written, parameterizable $W{\times}W$ PE-array accelerator executes this integer model bit-exactly; with $W{=}8$ it closes timing at 175.4~MHz out of context and at 142.86~MHz in a complete Zynq-7000 system, using 3,117 LUTs, 104 DSP slices and 14 block RAMs. One inference takes 142.2~$\mu$s, and measured on the ZedBoard, 161.9~$\mu$s including input transfer, 60.5$\times$ faster than the same model on the on-chip ARM Cortex-A9. Future work includes a jamming/spoofing/interference dataset with jammer-to-signal-ratio sweeps and classical detector baselines, board-level power measurement, and streaming input over AXI4-Stream.

\begin{thebibliography}{00}

\bibitem{b1} A. C. S. Thakur and M. Imtiaz, ``Deep learning for automatic modulation classification: A review,'' \emph{Electronics}, vol. 15, no. 10, p. 2163, May 2026.

\bibitem{b2} H. Xu, Y. Liang, R. Xie, and Y. Kong, ``The deep learning evolution in wireless physical layer communications: Applications, challenges, and evolutionary directions,'' \emph{Sensors}, vol. 26, no. 11, p. 3609, Jun. 2026.

\bibitem{b3} DeepSig Inc., ``RadioML 2018.01A dataset,'' 2018. [Online]. Available: https://www.deepsig.ai/datasets

\bibitem{b4} S. Ansari et al., ``Attention-enhanced hybrid automatic modulation classification for advanced wireless communication systems: A deep learning-transformer framework,'' \emph{IEEE Access}, vol. 13, 2025.

\bibitem{b5} F. Zhu et al., ``Robust deep learning-based physical layer communications: Strategies and approaches,'' \emph{IEEE Netw.}, 2025.

\bibitem{b6} S. Chatterjee, S. Ghosh, T. Ghosh, and H. Rahaman, ``Architectural design and performance analysis of FPGA based AI accelerators: A comprehensive review,'' arXiv preprint arXiv:2603.08740, Feb. 2026.

\bibitem{b7} X. Chen, X. Tian, Q. Zheng, and W. Wang, ``Attention mechanism enabled wireless modulation recognition network,'' \emph{Eng. Lett.}, vol. 33, no. 4, Apr. 2025.

\bibitem{b8} E. Chang, D. Cheng, and F. T. Luk, ``DeepRadio: Real-time FPGA deep learning for RF signal classification in electronic warfare,'' in \emph{Proc. IEEE Mil. Commun. Conf. (MILCOM)}, Norfolk, VA, Nov. 2019, pp. 1--6.

\bibitem{b9} E. Chang \emph{et al.}, ``Real-time modulation classification on Xilinx RFSoC using deep neural networks,'' in \emph{Proc. IEEE Int. Symp. Circuits Syst. (ISCAS)}, Seville, Spain, May 2020, pp. 1--5.

\bibitem{b10} D. J. M. Moss, D. Boland, P. Pourbeik, and P. H. W. Leong, ``Real-time FPGA-based anomaly detection for radio frequency signals,'' in \emph{Proc. IEEE Int. Symp. Circuits Syst. (ISCAS)}, Florence, Italy, May 2018, pp. 1--5.

\bibitem{b11} T. J. O'Shea, T. Roy, and T. C. Clancy, ``Over-the-air deep learning based radio signal classification,'' \emph{IEEE J. Sel. Topics Signal Process.}, vol. 12, no. 1, pp. 168--179, Feb. 2018.

\bibitem{b12} J. Duarte \emph{et al.}, ``Fast inference of deep neural networks in FPGAs for particle physics,'' \emph{J. Instrum.}, vol. 13, no. 07, p. P07027, Jul. 2018.

\bibitem{b13} X. Wang, Y. Zhao, and Z. Huang, ``A survey of deep transfer learning in automatic modulation classification,'' \emph{IEEE Trans. Cogn. Commun. Netw.}, vol. 11, no. 3, pp. 1357-1381, June 2025.

\bibitem{b14} L. Zhang, T. Zhou, J. Yang, Y. Li, Z. Zhang, X. Hu, and Y. Peng, ``A real-time signal measurement system using FPGA-based deep learning accelerators and microwave photonic,'' \emph{Remote Sens.}, vol. 16, no. 23, p. 4358, Nov. 2024.

\bibitem{b15} K. He, X. Zhang, S. Ren, and J. Sun, ``Deep residual learning for image recognition,'' in \emph{Proc. IEEE Conf. Comput. Vis. Pattern Recognit. (CVPR)}, Las Vegas, NV, Jun. 2016, pp. 770--778.

\bibitem{b16} Y. LeCun, L. Bottou, Y. Bengio, and P. Haffner, ``Gradient-based learning applied to document recognition,'' \emph{Proc. IEEE}, vol. 86, no. 11, pp. 2278--2324, Nov. 1998.

\bibitem{b17} S. Ioffe and C. Szegedy, ``Batch normalization: Accelerating deep network training by reducing internal covariate shift,'' in \emph{Proc. 32nd Int. Conf. Mach. Learn. (ICML)}, Lille, France, Jul. 2015, pp. 448--456.

\bibitem{b18} B. Jacob \emph{et al.}, ``Quantization and training of neural networks for efficient integer-arithmetic-only inference,'' in \emph{Proc. IEEE CVPR}, Salt Lake City, UT, Jun. 2018, pp. 2704--2713.

\bibitem{b19} D. P. Kingma and J. Ba, ``Adam: A method for stochastic optimization,'' in \emph{Proc. Int. Conf. Learn. Representations (ICLR)}, San Diego, CA, USA, May 2015.

\bibitem{b20} Y. Umuroglu et al., ``FINN: A framework for fast, scalable binarized neural network inference,'' in \emph{Proc. ACM/SIGDA Int. Symp. Field-Program. Gate Arrays (FPGA)}, Monterey, CA, Feb. 2017, pp. 65--74.

\end{thebibliography}

\end{document}
